\section{Реализация выборок и соединений (DQL)}
\label{sec:dql}

В приложении представлены два ключевых механизма: \texttt{SELECT} и \texttt{INNER JOIN}.

\subsection{Параметризованная выборка: POST /query/select}

Конструктор выборки \texttt{/select} позволяет пользователю интерактивно настраивать три компонента SQL-запроса:
\begin{enumerate}
    \item проекция колонок;
    \item предикаты фильтрации WHERE;
    \item сортировка ORDER BY.
\end{enumerate}

Фрагмент формы настройки выборки представлен в листинге~\ref{lst:select-form}.

\begin{listing}[H]
    \htmlfile[firstline=19, lastline=39, firstnumber=19]{../templates/select.html}
    \caption{Интерактивный конструктор выборки в \texttt{templates/select.html}}
    \label{lst:select-form}
\end{listing}

Обработка запроса выполняется в функции \texttt{SelectHandler} (листинг~\ref{lst:select-handler}).

\begin{listing}[H]
    \gofile[firstline=101, lastline=125, firstnumber=101]{../handlers_query.go}
    \caption{Сборка условий фильтрации, сортировки и запроса в \texttt{handlers\_query.go}}
    \label{lst:select-handler}
\end{listing}

Результат запроса считывается методом \texttt{pool.Query(querySQL)}, оборачивается в структуру \texttt{TableData} и отображается пользователю.

\subsection{Внутреннее соединение таблиц: POST /query/join}

Соединение таблиц по маршруту \texttt{/join} демонстрирует операцию внутреннего соединения (\texttt{INNER JOIN}). В качестве демонстрационной схемы используются таблицы \texttt{lab2\_employee\_ilinov} и \texttt{departments}, связанные отношением внешнего ключа: столбец \texttt{department\_id} таблицы сотрудников ссылается на первичный ключ \texttt{id} таблицы отделов.

Обработчик \texttt{JoinHandler} выполняет валидацию и собирает итоговый запрос (листинг~\ref{lst:join-handler}).

\begin{listing}[H]
    \gofile[firstline=235, lastline=249, firstnumber=235]{../handlers_query.go}
    \caption{Сборка и исполнение команды \texttt{INNER JOIN} в \texttt{handlers\_query.go}}
    \label{lst:join-handler}
\end{listing}